
\section*{Solutions to Exercises from Section 3}
\begin{solution}\label{sol:3.1}\ref{ex:3.1}
    
    $$
    \begin{aligned}
        &(1)~ \{ 1,2,3,4 \}=\\
        &\varnothing, \{ 1 \},\{ 2 \},\{ 3 \},\{ 4 \},\\
        &\{ 1,2 \},\{ 1,3 \},\{ 1,4 \},\{ 2,3 \},\\
        &\{ 2,4 \}, \{ 3,4 \},\{ 1,2,3 \},\{ 1,2,4 \},\\
        &\{ 1,3,4 \}, \{ 2,3,4 \}, \{ 1,2,3,4 \}\\\\
        &(2)~\{ 1,2,\varnothing \}=\\
        &\varnothing, \{ 1 \},\{ 2\},\\
        &\{ \varnothing \}, \{ 1,2 \},\{ 1,\varnothing \},\\
        &\{ 2,\varnothing \}, \{ 1,2,\varnothing \}\\\\
        &(3)~ \{ \{ \R \} \}=\varnothing, \{ \R \}\\\\
        &(4)~\varnothing = \varnothing\\\\
        &(5)~\{ \varnothing \}=\varnothing, \{ \varnothing \}\\\\
        &(6)~\{ \R,\Q,\N \}=\\
        &\varnothing, \{ \R \},\{ N \},\{ \Q \},\\
        &\{ \R,\N \},\{ \R,\Q \}, \{ \N,\Q \},\{ \R,\N,\Q \}\\\\
        &(7)~\{ \R,\{ \N,\Q \} \}=\\
        &\varnothing, \{ \R \},\{\{ \N,\Q \}\}, \{ \R,\{ \N,\Q \} \}\\\\
        &(8)~\{ \{ 0,1 \},\{ 0,1,\{ 2 \} \},\{ 0 \} \}=\\
        &\varnothing, \{\{ 0,1 \}\}, \{\{ 0,1,\{ 2 \} \}\},\{\{ 0 \}\},\\
        &\{ \{ 0,1 \},\{ 0,1,\{ 2 \} \} \}, \{ \{ 0,1 \}, \{ 0 \} \},\\
        &\{ \{ 0,1,\{ 2 \} \},\{ 0 \} \}, \{ \{ 0,1 \},\{ 0,1,\{ 2 \} \},\{ 0 \} \}.
    \end{aligned}
    $$
\end{solution}

\begin{solution}\label{sol:3.2}\ref{ex:3.2}
    $$
    \begin{aligned}
    &(1)~\{ X:X\subseteq \{ 3,2,a \}\text{ and } \lvert X \rvert=2 \}=\{ \{ 3,2 \},\{ 3,a \},\{ 2,a \} \}\\\\
    &(2)~\{ X\subseteq\N:\lvert X \rvert\le 1 \}=\{ \varnothing \}\cup \{ \{ n \}:n\in\N \}\\\\
    &(3)~\{ X:X\subseteq \{ 3,2,a \}\text{ and }\lvert X \rvert=4 \}=\varnothing\\\\
    &(4)~\{ X:X\subseteq \{ 3,2,a \}\text{ and }\lvert X \rvert=1 \}=\{ \{ 3 \},\{ 2 \},\{ a \} \}\\\\
    \end{aligned}
    $$ 
   \end{solution}

\begin{solution}\label{sol:3.3}\ref{ex:3.3}
   
   $$
   (1)~\R^3\subseteq\R^3
   $$

   $(1)$ is true, for the fact that every set is a subset of itself. 
   
   $$
   (2)~\R^2\subseteq\R^3
   $$
   $(2)$ is tricky, strictly speaking it is false. However, if you have any experience with Linear Algebra you might consider it true. In Linear Algebra we abuse the notation and say that $\R^3$ is an extension of $\R^2$, addiing the third dimension to a two dimensional plane. In set theory however, we are a bit more strict. The above is false, since 
   $$
   \R^2=\R\times\R=\{ (x,y):x\in \R\land y\in\R \},
   $$
   And 
   $$
   \R^3=\R\times\R\times\R=\{ (x,y,z):x\in\R\land y\in\R \land z\in\R \}.
   $$
   So, the first \emph{object} takes in pairs, and the second object takes in triples. 

   $$
   (3)~\{ (x,y)\in\R^2:x-1=0 \}\subseteq \{ (x,y)\in\R^2:x^2-x=0 \}.
   $$
   Consider a small example here, the set contains pairs $(x,y)$, the first element of the pair has to be $1$, so $\{ (1,y):y\in\R \}$. In the second set we have $x^2-x=0$, solving this gives us $x=0,\lor x=1$, so we have the pairs $\{ (0,y),(1,y):y\in\R \}$. Then it follows that the example is true. 
   $$
   (4)~\{ (x,y)\in\R^2:x^2-x=0 \}\subseteq \{ (x,y)\in\R^2:x-1=0 \}.
   $$
   This example is false, in the LHS, we can have $0^2-0=0$, whereas in the RHS we cannot, unless $x=1$.  
\end{solution}